% !TeX root = ../navier-stokes-manuscript.tex
% ============================================================
% 2.7 Compatibility and the Terminal Quantifier
% ============================================================

\subsection{Compatibility and the Terminal Quantifier}\label{sub:CompatibilityAndTheTerminalQuantifier}

Begin with two concrete finite views on the same strict interval \((0,T)\).
The larger record \(\cR_{2,3}(u,p;T)\) contains
\[
  u,
  \ \partial_tu,
  \ \partial_t^2u
  \in
  L^2(0,T;H^4)
\]
and
\[
  \partial_tu,
  \ \partial_t^2u,
  \ \partial_t^3u
  \in
  L^2(0,T;H^2).
\]
To recover the smaller record \(\cR_{1,2}(u,p;T)\), discard the highest temporal entries, use \(H^4\hookrightarrow H^3\) on the upper row, and use \(H^2\hookrightarrow H^1\) on the lower row.
The surviving coordinates are exactly
\[
  u,
  \ \partial_tu
  \in
  L^2(0,T;H^3)
\]
and
\[
  \partial_tu,
  \ \partial_t^2u
  \in
  L^2(0,T;H^1).
\]
The operation changes the resolution of the view while leaving every surviving derivative unchanged.
The pressure coordinates restrict through the same differentiated Poisson equations, so pressure remains attached to the smaller record.

The general restriction has three independent parts.
Let
\[
  N'\leq N,
  \qquad
  M'\leq M,
  \qquad
  0<T'\leq T<T_*.
\]
Define
\[
  \rho_{N',M';T'}^{N,M;T}:
  \cR_{N,M}(u,p;T)
  \longrightarrow
  \cR_{N',M'}(u,p;T')
\]
by forgetting temporal rows above \(N'\), applying the continuous Sobolev embeddings down to the spatial heights determined by \(M'\), and restricting time from \((0,T)\) to \((0,T')\).
Every shared velocity coordinate remains the same distributional derivative \(\partial_t^m\partial_x^\alpha u\).
Every shared pressure coordinate remains the same derivative \(\partial_t^m\partial_x^\alpha p\).
Every retained coordinate continues to satisfy the corresponding row of the pressure-complete mixed recurrence.

Restriction in stages agrees with restriction in one step.
If
\[
  N''\leq N'\leq N,
  \qquad
  M''\leq M'\leq M,
  \qquad
  T''\leq T'\leq T,
\]
then
\[
  \rho_{N'',M'';T''}^{N',M';T'}
  \circ
  \rho_{N',M';T'}^{N,M;T}
  =
  \rho_{N'',M'';T''}^{N,M;T}.
\]
This composition law is compatibility.
A higher-resolution finite view always returns the same lower-resolution view, independent of how many intermediate cutoffs are used.
The compatible family can support the coordinatewise all-orders intersection discovered above.

Compatibility settles agreement among the coordinates.
An alleged finite maximal time introduces a separate question about how far their bounds extend.
Fix a strict \(T<T_*\).
Classical smoothness on the compact interval \([0,T]\) gives, for every finite \(n\) and \(q\),
\[
  \sup_{0\leq t\leq T}
  \norm{\partial_t^n u(t)}_{H^q}
  <\infty.
\]
Integrating this fixed-cutoff bound gives
\[
  \norm{\partial_t^n u}_{L^2(0,T;H^q)}^2
  \leq
  T
  \sup_{0\leq t\leq T}
  \norm{\partial_t^n u(t)}_{H^q}^2
  <\infty.
\]
Hence every finite rectangle is finite after a strict cutoff has been fixed.

It is useful to write the dependence of the bound explicitly.
Interior classicality gives
\[
  \forall\,T<T_*
  \ \forall\,N,M\in\Nzero
  \ \exists\,C_{T,N,M}<\infty:
  \qquad
  \mathfrak R_{N,M}(u;T)
  \leq
  C_{T,N,M}.
\]
The constant is allowed to depend on the right endpoint \(T\).
This statement permits one fixed rectangle to grow without bound along a sequence \(T_j\uparrow T_*\), even though its value is finite at every earlier cutoff.

The whole-terminal statement moves the existential constant outside the time cutoff:
\[
  \forall\,N,M\in\Nzero
  \ \exists\,C_{N,M}(u_0,\nu,T_*)<\infty
  \ \forall\,T<T_*:
  \qquad
  \mathfrak R_{N,M}(u;T)
  \leq
  C_{N,M}(u_0,\nu,T_*).
\]
Equivalently, for every fixed finite pair \((N,M)\),
\[
  \sup_{0<T<T_*}
  \mathfrak R_{N,M}(u;T)
  \leq
  C_{N,M}(u_0,\nu,T_*)
  <\infty.
\]
The constant may change when \(N\) or \(M\) changes.
Its essential uniformity is across every strict time cutoff for the selected rectangle.

The distinction lies in the placement of the bound, not in the order of two universal quantifiers.
The strict-subinterval statement allows \(C_{T,N,M}\) to deteriorate as \(T\) approaches \(T_*\).
The whole-terminal statement fixes \((N,M)\), supplies one datum- and viscosity-dependent constant, and carries that coordinate all the way to the terminal face.

The datum-generated whole-terminal rectangle theorem supplies precisely this statement for every finite rectangle.
The pressure-complete mixed jet, read through every finite adjacent Sobolev rectangle with all overlapping coordinates compatible, is the object controlled by that theorem.
Once each fixed rectangle has its cutoff-independent bound, no mixed coordinate can escape by moving the right endpoint of integration toward \(T_*\).
The compatible family is then available across the whole alleged terminal interval for the formal construction that follows.
